Large language models (LLMs) are increasingly used for software engineering tasks such as vulnerability detection, but standard LLM inference treats each input independently and does not preserve past failures. As a result, a model can repeatedly make similar false-positive or false-negative predictions. This problem is important for CWE-based vulnerability detection because false positives waste developer effort, while false negatives can leave real security bugs undiscovered. We propose \system, a database-backed mistake-memory framework that records incorrect LLM predictions and retrieves similar past mistakes during future inference. Instead of modifying the base model, \system adds a persistent SQLite mistake database and a retrieval-augmented prompting layer. Each stored mistake includes the code sample, ground-truth label, predicted label, CWE information, raw model output, error type, and retrieval metadata. We evaluate \system using Qwen2.5-Coder-7B-Instruct on CASTLE-C250 and a sampled Juliet C/C++ Test Suite subset with CWE categories overlapping CASTLE. CASTLE-C250 is a hand-crafted CWE benchmark with 250 C micro-benchmark programs covering 25 common CWEs, while Juliet is a larger C/C++ test suite organized under 118 CWE categories. Our results show mixed behavior on CASTLE: \system reduces false positives but can increase false negatives, making the model more conservative. However, on the external Juliet subset, the CASTLE-built mistake database improves accuracy from 68.67\% to 82.40\% and macro F1 from 66.23\% to 82.18\%. We further add a supplemental full-C250 compact-RAG experiment using OpenRouter's \texttt{openai/gpt-oss-120b:free}, where \system improves F1 from 0.616 to 0.709 and CWE correctness from 0.433 to 0.520. These results suggest that persistent mistake memory is a promising direction for improving LLM vulnerability detection, while also showing the need for selective retrieval, gating, and prompt balancing.
